% =======================================================
% 2. RELATED WORK
% =======================================================
\section{Related Work}
\label{sec:related}

\paragraph{Static and polyhedral analysis} Classical tests (GCD, Banerjee, the Omega test) decide whether two affine array references can touch the same location in different iterations \cite{banerjee1988,pugh1992}. LLVM's DependenceAnalysis implements a family of such tests on top of ScalarEvolution \cite{pop2005} and alias analysis and returns direction and distance vectors per loop level \cite{lattner2004,llvmda}. The polyhedral model generalizes dependence analysis to whole loop nests and enables tiling and interchange \cite{feautrier1991,bondhugula2008,grosser2012}, but it requires static control parts with affine bounds and subscripts. Production compilers expose auto-parallelization as an option (GCC \Code{-ftree-parallelize-loops}, Polly \Code{-polly-parallel}); both are sound and both reject loops whenever aliasing or subscript structure cannot be resolved at compile time. Dynamic tools such as DiscoPoP \cite{li2016discopop} observe dependences at run time and depend on input coverage.

\paragraph{Learning-based parallelism detection} PragFormer \cite{harel2023pragformer} and OMPify \cite{kadosh2023ompify} train transformers to predict whether a loop should carry an OpenMP directive and which clauses it needs; OMPar \cite{kadosh2024ompar} pairs such a classifier with a code generator. Graph-based models learn from flow-aware program representations \cite{chen2023graph2par}, and AutoParLLM \cite{mahmud2025autoparllm} uses a GNN to build the context of an LLM prompt. ParEval \cite{nichols2024pareval} found that general-purpose LLMs generate correct parallel code far less often than correct serial code, and LLM-Vectorizer \cite{taneja2025llmvectorizer} pairs LLM-generated vector code with formal translation validation. Concerns about the correctness of model-generated code in general were raised early \cite{pearce2022asleep}.

\paragraph{Benchmarks and ground truth} DataRaceBench \cite{liao2017dataracebench} provides micro-benchmarks with known race status for evaluating race detectors, i.e.\ it labels \emph{existing} directives. Learning-based studies usually mine labels from repositories (a loop is ``parallel'' if a developer annotated it), which does not establish that unannotated loops are serial. PolyBench \cite{pouchet2012} provides affine kernels without legality labels. Our benchmarks label \emph{sequential} loops under an explicit legality definition and check every label by execution.

\paragraph{Positioning} Retrieval-augmented generation and agentic workflows can supply documentation to a model, but they do not supply the one fact that decides legality---whether a dependence exists for this loop. \SystemName{} differs from the learning-based systems above in that the model is never the final authority on a loop the compiler can decide, and from the compiler-based systems in that undecided loops are not automatically discarded. The closest idea is AutoParLLM's use of program analysis to build prompt context; in \SystemName{} the analysis is a sound oracle with a veto, and the context it provides is the reason for its own uncertainty.
